% Section 2, Related Work.

\section{Related Work}
\label{sec:related}

SWE-agent~\citep{sweagent} introduced the agent-computer-interface paradigm
for autonomous issue resolution in real repositories.
Agentless~\citep{agentless} showed that a fixed pipeline, with no agent
autonomy at all, can achieve competitive resolve rates.
SWE-smith~\citep{swesmith} addresses data scarcity by synthesizing
large-scale training corpora for such agents, while
SWE-Gym~\citep{swegym} and R2E-Gym~\citep{r2egym} supply executable
training environments built from real repositories.
The community evaluates these systems on
SWE-bench~\citep{swebench} and its harder successor
SWE-bench Pro~\citep{swebenchpro}.

A parallel line of work trains small models specifically for code
localization.
SWE-Fixer~\citep{swefixer} trains a 7B retriever whose output feeds one
fixed larger editor.
LocAgent~\citep{locagent} and SweRank~\citep{swerank} similarly train
compact models to identify fault locations within a repository.
These systems produce exactly the kind of evidence a router could consume,
yet none of them routes: the localizer's output terminates in a single,
predetermined consumer.

LLM routing has been studied along two axes.
Cost-quality routers such as RouteLLM~\citep{routellm},
FrugalGPT~\citep{frugalgpt}, and Hybrid LLM~\citep{hybridllm} learn to
dispatch queries to cheaper or stronger models based on predicted difficulty.
Profile-based selectors take a complementary approach, building per-model
signatures from benchmark outcomes or learned
embeddings and matching incoming queries against
them~\citep{shnitzer,embedllm,iclrouter}.
Both families share a structural property: they decide from the task prompt,
and optionally from model profiles, before any system has engaged with the
concrete problem instance.

Several concurrent efforts address routing in code-generation settings.
SWE-Router~\citep{swerouter} probes each issue with a weak model and then
escalates between exactly two models; on escalation the strong model restarts
from scratch, so the probe's work is not consumed.
Its theoretical analysis of trajectory-conditioned routing nonetheless
supports the design direction we pursue.
CodeRescue~\citep{coderescue} routes among recovery actions (reflect,
replan, or escalate) inside a single agent's trajectory rather than among
fixer models.
TRACE-Router~\citep{tracerouter} formulates model selection as a contextual
bandit over an N-way pool, though it is evaluated outside software
engineering.
Self-play SWE-RL~\citep{selfplayswerl} trains solvers via self-play and
evaluates on SWE-bench Pro but does not route.
As Table~\ref{tab:capability-matrix} summarizes, \sysname is distinct in
combining a trained searcher, a handoff the fixer actually consumes (after
verification), an N-way fixer pool where adding a new fixer requires no
retraining, and routing features drawn from the searcher's hidden states.

\begin{table}[!tbp]
\centering
\caption{\textbf{Capability comparison of \sysname with related routing and
localization systems.}
Columns: \emph{Searcher}, a model trained to scout the repository before
any fix is attempted;
\emph{Handoff}, the searcher's work product is consumed by the downstream
fixer rather than discarded;
\emph{$N$-way}, selection over a pool of more than two fixers;
\emph{Onboard}, adding a new fixer requires no retraining of any learned
component;
\emph{States}, routing features include the searcher's internal hidden
states.
Concurrent code-routing work escalates between a fixed pair and discards
the cheap model's evidence on escalation;
profile-based routers select from the task prompt alone, before any
exploration;
trained small localizers produce evidence a router could use but perform no
routing.
$\checkmark^{\dagger}$: SWE-Fixer's retriever feeds one fixed larger editor,
while LocAgent and SweRank stop at localization.
Rows: SWE-Router~\citep{swerouter}, CodeRescue~\citep{coderescue},
TRACE-Router~\citep{tracerouter},
profile routers~\citep{embedllm,iclrouter,shnitzer},
small localizers~\citep{swefixer,locagent,swerank}.}
\label{tab:capability-matrix}
% ===========================================================================
% tab-capability-matrix.tex — related-work capability matrix (Section 2).
% BODY ONLY (tabular). Wrapper/caption/label live in the section file.
%
% PROVENANCE — every cell.
%   Ground truth for the three code-routing competitors = the scoop-check
%   entries in internal records (each paper was READ
%   IN FULL before the verdict was written).
%
%   SWE-Router (arXiv 2607.00053, ICML26 DL4Code workshop) — internal
%   scoop-check entry:
%     trained searcher   = NO. Their probe is the WEAK FIXER run for K=3
%                          turns; the trained 7B is a value head reading the
%                          partial trajectory, not a searcher. Entry lists
%                          "trained-7B searcher" under "Untouched by them".
%     consumed handoff   = NO. "on escalation the strong model RESTARTS FROM
%                          SCRATCH (probe discarded) ... their probe is
%                          thrown away — citable gap".
%     N-way pool         = NO. "binary continue-vs-escalate", "per-pair
%                          value head".
%     zero-retrain onb.  = NO. "fresh rollouts + retraining per model".
%     hidden states      = NO. Entry: "hidden states (text-only)".
%
%   CodeRescue (arXiv 2607.19338, preprint) — internal scoop-check entry:
%     trained searcher   = NO. "a 4B router picks among 3 fixed ACTIONS
%                          (reflect/replan/escalate)"; entry lists "no
%                          searcher/handoff/resumes/hidden-states".
%     consumed handoff   = NO (same clause).
%     N-way pool         = NO. "action-routing not model-routing, binary
%                          cheap/strong pair not N-way pool"; also
%                          "function-level not repo-level".
%     zero-retrain onb.  = NO. "their fixed 3-way head structurally cannot
%                          admit a new model without retraining".
%     hidden states      = NO (same "no ... hidden-states" clause).
%
%   TRACE-Router (arXiv 2607.22465) — internal records characterises it
%   only as "task-level bandit routing, non-SWE, low risk", and the
%   competitive-map summary states that NONE of the three has "a
%   trained searcher, consumed handoff, resume N-way pool, zero-retraining
%   onboarding, handoff-conditioned labels, hidden-state features, or a Pro
%   eval". That licenses four of its five cells directly.
%     N-way pool cell WEB-VERIFIED: their framework is defined
%     over "a pool of K candidate LLM backends" (arbitrary K) and §4.6
%     demonstrates a 4-arm experiment, though headline results use pairs.
%     Check mark stands (arxiv.org/abs/2607.22465).
%
%   Profile-based routers (EmbedLLM / ICL-Router / Shnitzer, one collapsed
%   row) — internal records (positioning, scoop-checked):
%     "Novelty wedge = multi-fixer selection AFTER a dedicated reconnaissance
%      phase, via an explicit handoff. Cite & differentiate: ... EmbedLLM/
%      ICL-Router/Shnitzer (profiles-from-outcomes concept — NOT ours to
%      claim)".
%     trained searcher   = NO  (the wedge sentence: they route without any
%                          reconnaissance phase; nothing is trained to search)
%     consumed handoff   = NO  (same sentence — no handoff exists to consume)
%     N-way pool         = YES (the profiles-from-outcomes concept is a
%                          many-model profile bank; internal records credits the concept
%                          to them explicitly)
%     zero-retrain onb.  = YES (same concept: a new model joins by having its
%                          outcome profile built; internal records describes our own
%                          resume mechanism as the same idea and calls the
%                          concept "NOT ours to claim")
%     hidden states      = NO
%     hidden-states cell WEB-VERIFIED: EmbedLLM routes on a
%     correctness-matrix factorization + question text embeddings
%     (arxiv.org/abs/2410.02223); ICL-Router on profiling-derived in-context
%     vectors (arxiv.org/abs/2510.09719); Shnitzer on input-text embeddings
%     over benchmark outcomes (arxiv.org/abs/2309.15789, corroborated from
%     abstract; verbatim methods quote not pulled -- recheck only if a
%     reviewer contests). None access the routed models' activations; the
%     cross stands.
%
%   Small localizers (SWE-Fixer / LocAgent / SweRank, one collapsed row) —
%   internal records: "SWE-Fixer/LocAgent/Kimi-Dev (small localizers/
%   test-writers, no routing)"; internal records names SweRank-7B as the published
%   small-model localization baseline row.
%     trained searcher   = YES (that is exactly what they are — trained small
%                          localizers; internal records cites LocAgent as "closest
%                          task, LoRA on Qwen2.5-Coder")
%     consumed handoff   = MIXED, marked with a dagger. internal records
%                          (exam-menu sweep) records "SWE-Fixer (7B retr+72B
%                          editor) 32.8" — its retrieval output IS consumed
%                          by a fixed larger editor, but LocAgent and SweRank
%                          stop at localization. The dagger and the caption
%                          carry the distinction rather than hiding it.
%     N-way pool         = NO  ("no routing")
%     zero-retrain onb.  = NO  (no pool to onboard into)
%     hidden states      = NO  (retrieval/agent text pipelines)
%
%   \sysname row — all five are ours by construction:
%     trained searcher   = internal records (\modelname performs the
%                          search phase, SFT-trained)
%     consumed handoff   = harness rules (verify-then-strip, then the
%                          handoff is injected into the fixer prompt)
%     N-way pool         = four fixers, internal pool ruling plus the
%                          router spec's "pool" field
%     zero-retrain onb.  = internal records: "onboard new fixer from 25-50 outcomes,
%                          zero retraining"
%     hidden states      = the deployed head blends a task-text head with a
%                          \modelname pre-decode hidden-state head
%
%   WORDING RULE APPLIED: the
%   onboarding column is phrased "adding a new FIXER requires no retraining",
%   never a bare "without retraining" (that phrase collides with CodeRescue's
%   budget-knob usage).
% ===========================================================================
\begin{tabular}{@{}lccccc@{}}
\toprule
System & Searcher & Handoff & $N$-way & Onboard & States \\
\midrule
SWE-Router      & $\times$     & $\times$ & $\times$     & $\times$     & $\times$ \\
CodeRescue      & $\times$     & $\times$ & $\times$     & $\times$     & $\times$ \\
TRACE-Router    & $\times$     & $\times$ & $\checkmark$ & $\times$     & $\times$ \\
Profile routers & $\times$     & $\times$ & $\checkmark$ & $\checkmark$ & $\times$ \\
Small localizers & $\checkmark$ & $\checkmark^{\dagger}$ & $\times$ & $\times$ & $\times$ \\
\midrule
\textbf{\sysname} & $\checkmark$ & $\checkmark$ & $\checkmark$ & $\checkmark$ & $\checkmark$ \\
\bottomrule
\end{tabular}

\end{table}
\FloatBarrier
